\documentclass[12pt]{article}

\usepackage[utf8]{inputenc}
\usepackage[spanish]{babel}
\usepackage{amsmath, amssymb}
\usepackage{geometry}
\usepackage{enumitem}

\geometry{margin=2.5cm}

\title{Solución Analítica de los Modelos Dinámicos del Amor}
\author{}
\date{}

\begin{document}

\maketitle

\section{Modelo Lineal: La pareja Romeo y Julieta}

Para representar la evolución del amor entre la pareja de Romeo y Julieta con el paso del tiempo, se considera
el siguiente sistema lineal:

\begin{equation}
\begin{cases}
\frac{dR}{dt} = aR + bJ,\\
\frac{dJ}{dt} = cR + dJ.
\end{cases}
\end{equation}

Cada elemento representa lo siguiente:
	\begin{enumerate}[label=-]
	\item $R(t)$: sentimientos de Romeo hacia Julieta.
	\item $J(t)$: sentimientos de Julieta hacia Romeo.
	\item $a$: influencia de los sentimientos propios de Romeo.
	\item $b$: influencia de los sentimientos de Julieta sobre Romeo.
	\item $c$: influencia de los sentimientos de Romeo sobre Julieta.
	\item $d$: influencia de los sentimientos propios de Julieta.
	\end{enumerate}

Los parámetros $a$, $b$, $c$ y $d$ pueden tomar valores negativos, positivos o nulos.

$ $

Escrito en forma matricial, queda de la siguiente forma:

\begin{equation}
Y'=AY,
\qquad
A=
\begin{pmatrix}
a & b\\
c & d
\end{pmatrix},
\qquad
Y=
\begin{pmatrix}
R\\
J
\end{pmatrix}.
\end{equation}

\subsection{Solución de Equilibrio}
El origen $(0; 0)$ representa el único punto de equilibrio del modelo, por tanto $R(t) = 0$ y $J(t) = 0$ constituyen
las soluciones triviales del sistema.

\subsection{Hallar los valores propios}

\begin{equation}
det(A - \lambda I)= 0
\end{equation}

\begin{equation}
    \begin{vmatrix}
    a - \lambda & b\\
    c & d - \lambda
    \end{vmatrix} = 0
\end{equation}

\begin{equation}
\lambda^2-(a+d)\lambda+(ad-bc)=0
\end{equation}

A partir de $(5)$ obtenemos el discriminante:

\begin{equation}
\Delta=(a+d)^2-4(ad-bc)
\end{equation}

Luego, los valores propios serán:

\begin{equation}
\lambda_{1, 2} = \frac{a + d \pm \sqrt{(a+d)^2-4(ad-bc)}}{2}
\end{equation}

Debido a que los parámetros pueden ser negativos, positivos o nulos, puede suceder en consecuencia que el discriminante sea negativo,
positivo o nulo. De aquí se derivan varios casos en los que varía la forma de la solución general.

\subsection{Caso 1: $\Delta > 0 \Rightarrow$ Valores propios reales distintos}

Si $\lambda_1 \neq \lambda_2$, buscamos los vectores propios
$K_1$ y $K_2$ asociados respectivamente, los cuales serán l.i.

La solución general es:

\begin{equation}
Y(t)=c_1e^{\lambda_1t}K_1
+c_2e^{\lambda_2t}K_2.
\end{equation}

\subsection{Caso 2: $\Delta = 0 \Rightarrow$ Valor propio repetido}

Si $\lambda_1=\lambda_2=\lambda$, existen dos posibilidades.

\textbf{Dos vectores propios (l.i.):}

\begin{equation}
Y(t)=c_1e^{\lambda t}K_1 + c_2e^{\lambda t}K_2
\end{equation}

\textbf{Un solo vector propio:}

Buscamos un vector propio generalizado $P$ que satisfaga:

\begin{equation}
(A-\lambda I)P=K.
\end{equation}

La solución general es:

\begin{equation}
Y(t)=c_1e^{\lambda t}K
+c_2e^{\lambda t}(tK+P).
\end{equation}

\subsection{Caso 3: $\Delta < 0 \Rightarrow$ Valores propios complejos}

Si los valores propios son $\lambda=\alpha\pm i\beta$, y un vector propio complejo es $K = B_1 + iB_2$,
obtenemos dos soluciones reales:

\begin{equation}
Y_1(t)=e^{\alpha t}
(B_1\cos(\beta t)-B_2\sin(\beta t)),
\end{equation}

\begin{equation}
Y_2(t)=e^{\alpha t}
(B_1\sin(\beta t)+B_2\cos(\beta t)).
\end{equation}

Por tanto, la solución general es:

\begin{equation}
Y(t)=c_1Y_1(t)+c_2Y_2(t).
\end{equation}

\section{Análisis de los casos particulares}

\subsection{Caso 1: Out of touch ($a = d = 0$)}

Este caso consiste en que Romeo y Julieta no responden a sus propios sentimientos, sino únicamente a los de la otra persona.
En este caso se obtienen los valores propios: $\lambda_{1, 2} = \pm \sqrt{bc}$.

\subsection{Caso 2: Fire and Ice ($c = -b$ y $a = -d$)}

Este caso consiste en que Romeo y Julieta son opuestos exactos en su manera de responder.
En este caso se obtienen los valores propios: $\lambda_{1, 2} = \pm \sqrt{a^2 - b^2}$.

\subsection{Caso 3: Peas in a Pod ($c = b$ y $a = d$)}

Este caso consiste en que Romeo y Julieta son idénticos en su manera de responder.
En este caso se obtienen los valores propios: $\lambda_{1, 2} = a \pm b$.

\end{document}